\documentclass[11pt]{article}

\usepackage[spanish, es-nodecimaldot]{babel}
\usepackage[utf8]{inputenc}
\usepackage[T1]{fontenc}
\usepackage{geometry}
\usepackage{amsmath,amssymb,amsthm,mathtools}
\usepackage{array}
\usepackage{enumitem}
\usepackage{hyperref}
\usepackage{xcolor}

\geometry{letterpaper, margin=2.5cm}
\setlength{\parindent}{0pt}
\setlength{\parskip}{0.6em}

\newcommand{\curso}{ADA - Análisis y Diseño de Algoritmos (2026-1)}
\newcommand{\tarea}{Tarea 3 - Algoritmos Voraces}
\newcommand{\nombre}{Miguel Angel Padilla Rosero}
\newcommand{\codigo}{8988878}
\newcommand{\profesores}{Carlos Ramírez y Camilo Rocha}
\newcommand{\fecha}{\today}
\newcommand{\seccion}{\textit{Por completar}}
\newcommand{\colaboradores}{\textit{Por completar}}
\newcommand{\fuentes}{Kleinberg \& Tardos, Erickson y Rocha}

\newcommand{\N}{\mathbb{N}}

\begin{document}

\begin{titlepage}
    \centering
    {\Large \textbf{\curso}\par}
    \vspace{1cm}
    {\huge \textbf{\tarea}\par}
    \vspace{1.3cm}

    \begin{tabular}{>{\bfseries}l l}
        Nombre: & \nombre \\
        Código: & \codigo \\
        Sección: & \seccion \\
        Profesores: & \profesores \\
        Fecha: & \fecha \\
        Colaboradores: & \colaboradores \\
        Fuentes: & \fuentes \\
    \end{tabular}

    \vfill
\end{titlepage}

\section*{Problemas conceptuales}

\section{Problema 1: Planificación de trabajos en El Goog}

\subsection*{Resumen del problema}

Se tienen \(N\) trabajos independientes. Cada trabajo \(J_i\) consta de dos etapas:
primero debe ser procesado en una única supercomputadora, y luego debe finalizarse
en uno de varios PCs. Si el trabajo \(J_i\) toma \(s_i\) unidades de tiempo en la
supercomputadora y \(f_i\) unidades de tiempo en un PC, entonces el objetivo es
determinar un orden para alimentar los trabajos a la supercomputadora que minimice
el tiempo en que \emph{todos} los trabajos han terminado completamente.

La observación clave es que los PCs trabajan en paralelo, mientras que la
supercomputadora trabaja de manera secuencial. Por ello, conviene liberar cuanto antes
aquellos trabajos cuyo tiempo de finalización en PC es más grande. A partir de esta
intuición se diseña un algoritmo voraz: ordenar los trabajos descendentemente por
\(f_i\), y luego recorrerlos en ese orden calculando el tiempo total de completación.

\subsection*{Especificación del problema}

\textbf{Entrada:}

Un arreglo
\[
A[0..N), \qquad N \geq 0,
\]
de parejas de números naturales, tal que
\[
A[i] = (s_i,f_i),
\]
donde:
\begin{itemize}[nosep]
    \item \(s_i\) es el tiempo de procesamiento del trabajo \(J_i\) en la supercomputadora;
    \item \(f_i\) es el tiempo de finalización del trabajo \(J_i\) en un PC.
\end{itemize}

\textbf{Salida:}

Un orden de los trabajos que minimiza el tiempo total de completación del sistema.

\subsection*{Función objetivo}

Sea \(\pi\) una permutación de \(\{0,\dots,N-1\}\). Si los trabajos se ejecutan en el orden
\[
J_{\pi(0)}, J_{\pi(1)}, \dots, J_{\pi(N-1)},
\]
entonces el tiempo en que el trabajo \(J_{\pi(i)}\) termina su paso por la supercomputadora es
\[
S_{\pi}(i)=\sum_{j=0}^{i} s_{\pi(j)}.
\]

Como luego ese trabajo aún necesita \(f_{\pi(i)}\) unidades de tiempo en un PC, su tiempo
de terminación total es
\[
C_{\pi}(i)=S_{\pi}(i)+f_{\pi(i)}.
\]

Por tanto, el tiempo de completación del orden \(\pi\) es
\[
T(\pi)=\max_{0\leq i<N} C_{\pi}(i)
      =\max_{0\leq i<N}\left(\sum_{j=0}^{i} s_{\pi(j)} + f_{\pi(i)}\right).
\]

Si \(N=0\), por convención
\[
T(\emptyset)=0.
\]

El objetivo es calcular un orden \(\pi^\ast\) tal que
\[
T(\pi^\ast)=\min_{\pi} T(\pi).
\]

\subsection*{Estrategia voraz}

La estrategia voraz es la siguiente:

\begin{quote}
Entre los trabajos aún no programados, escoger primero uno cuyo tiempo de
finalización en PC sea máximo.
\end{quote}

Equivalentemente, si se ordena todo el arreglo descendentemente por la segunda
componente \(f_i\), ese será el orden construido por la estrategia voraz.

La intuición es que un trabajo con gran valor de \(f_i\) puede convertirse en el cuello de
botella del sistema. Si se libera temprano hacia los PCs, su tiempo largo de finalización
transcurre en paralelo con la ejecución de otros trabajos en la supercomputadora.

\subsection*{Ejemplo}

Considérese el arreglo:
\[
\begin{array}{c|cccc}
i & 0 & 1 & 2 & 3\\
\hline
s_i & 2 & 5 & 1 & 4\\
f_i & 7 & 3 & 9 & 6
\end{array}
\]

El criterio voraz ordena los trabajos como
\[
J_2,\; J_0,\; J_3,\; J_1
\]
porque
\[
9 > 7 > 6 > 3.
\]

El recorrido produce:
\[
\begin{array}{c|cccc}
\text{trabajo} & J_2 & J_0 & J_3 & J_1\\
\hline
S(i) & 1 & 3 & 7 & 12\\
C(i)=S(i)+f_i & 10 & 10 & 13 & 15
\end{array}
\]

Luego, el tiempo total de completación es
\[
\max\{10,10,13,15\}=15.
\]

\subsection*{Pseudocódigo}

\[
\begin{array}{l}
\textsc{ElGoog}(A[0..N)):\\
\quad \text{ordenar } A \text{ descendentemente por la segunda componente} \\
\quad ac \gets 0 \\
\quad ans \gets 0 \\
\\
\quad \textbf{for } i = 0 \textbf{ to } N-1 \textbf{ do}\\
\quad \quad ac \gets ac + A[i][0] \\
\quad \quad ans \gets \max(ans,\, ac + A[i][1]) \\
\\
\quad \textbf{return } ans
\end{array}
\]

\subsection*{Correctitud}

\paragraph{Teorema de optimización local.}
Sea \(X\) un conjunto no vacío de trabajos y sea \(J_k \in X\) un trabajo tal que
\[
f_k=\max\{f_i : J_i\in X\}.
\]
Entonces existe un orden óptimo de los trabajos en \(X\) en el que \(J_k\) aparece en
la primera posición.

\paragraph{Demostración.}
Sea
\[
B=(b_0,b_1,\dots,b_{m-1})
\]
un orden óptimo de los trabajos en \(X\), donde \(m=|X|\). Si \(b_0=k\), no hay nada
que demostrar.

Supóngase entonces que \(b_r=k\) para algún \(r>0\). Considérese el orden
\[
B'=(k,b_0,b_1,\dots,b_{r-1},b_{r+1},\dots,b_{m-1}),
\]
es decir, el orden obtenido al mover \(J_k\) a la primera posición y desplazar una
posición a la derecha a los trabajos \(b_0,\dots,b_{r-1}\).

Sea
\[
u=\sum_{j=0}^{r-1} s_{b_j}.
\]

En el orden \(B\), el tiempo total de terminación del trabajo \(J_k\) es
\[
C_B(k)=u+s_k+f_k.
\]

En el orden \(B'\), ese mismo trabajo termina en
\[
C_{B'}(k)=s_k+f_k\leq C_B(k).
\]

Ahora sea \(b_j\) un trabajo con \(0\leq j<r\). En el orden \(B'\), su tiempo total de
terminación es
\[
C_{B'}(b_j)=s_k+\sum_{t=0}^{j} s_{b_t}+f_{b_j}.
\]
Como
\[
\sum_{t=0}^{j} s_{b_t}\leq u
\qquad\text{y}\qquad
f_{b_j}\leq f_k,
\]
se obtiene
\[
C_{B'}(b_j)\leq s_k+u+f_k=C_B(k).
\]

Finalmente, para todo trabajo que en \(B\) estaba después de la posición \(r\), la suma
de tiempos de supercomputadora de los trabajos anteriores no cambia, porque el conjunto
de trabajos antes de él sigue siendo el mismo. En consecuencia, su tiempo total de
terminación permanece igual.

Por tanto, todos los trabajos del bloque afectado en \(B'\) terminan a lo sumo en
\(C_B(k)\), y los demás no empeoran. Luego
\[
T(B')\leq T(B).
\]
Como \(B\) era óptimo, \(B'\) también es óptimo. En particular, existe un orden óptimo
con \(J_k\) en la primera posición. \(\square\)

\paragraph{Teorema de optimización global.}
Sea \(B[0..N)\) el arreglo obtenido al ordenar \(A[0..N)\) descendentemente por
la componente \(f_i\). Entonces \(B\) es un orden óptimo.

\paragraph{Demostración.}
Se probará por inducción sobre la longitud del prefijo que, para cada
\(t\), con \(0\leq t\leq N\), existe un orden óptimo cuyo prefijo de longitud \(t\) coincide
con \(B[0..t)\).

\textbf{Caso base:} \(t=0\). Es inmediato.

\textbf{Paso inductivo:} supóngase que existe un orden óptimo cuyo prefijo de longitud
\(t\) es
\[
B[0],B[1],\dots,B[t-1].
\]
Considérese el subproblema formado por los trabajos restantes. Por construcción del
arreglo \(B\), el trabajo \(B[t]\) tiene tiempo de finalización en PC máximo entre los
trabajos aún no fijados. Por el teorema de optimización local, existe un orden óptimo
de ese subproblema en el que \(B[t]\) aparece primero. Al anteponer el prefijo ya fijado,
se obtiene un orden óptimo global cuyo prefijo de longitud \(t+1\) coincide con
\[
B[0],B[1],\dots,B[t].
\]

Por inducción, existe un orden óptimo cuyo prefijo de longitud \(N\) es exactamente
\(B[0..N)\). Por lo tanto, \(B\) es óptimo. \(\square\)

\subsection*{Invariante del ciclo}

Una vez el arreglo está ordenado, el recorrido lineal calcula el valor de la función
objetivo. La invariante del ciclo es:

\begin{quote}
Al comienzo de la iteración \(i\), con \(0\leq i\leq N\),
\[
ac=\sum_{j=0}^{i-1} B[j][0]
\]
y
\[
ans=\max_{0\leq j<i}\left(\sum_{k=0}^{j} B[k][0] + B[j][1]\right),
\]
con la convención de que el máximo sobre el conjunto vacío vale \(0\).
\end{quote}

\paragraph{Demostración de la invariante.}
Para \(i=0\), antes de entrar al ciclo,
\[
ac=0
\qquad\text{y}\qquad
ans=0,
\]
así que la propiedad es cierta.

Supóngase que la propiedad es cierta al comienzo de la iteración \(i\), con \(0\leq i<N\).
La instrucción
\[
ac \gets ac + B[i][0]
\]
hace que
\[
ac=\sum_{j=0}^{i} B[j][0].
\]
Luego, la instrucción
\[
ans \gets \max(ans,\, ac+B[i][1])
\]
actualiza \(ans\) al máximo entre los tiempos ya calculados para \(B[0..i)\) y el tiempo
de terminación del trabajo \(B[i]\). Por tanto, al final de la iteración se cumple
\[
ans=\max_{0\leq j\leq i}\left(\sum_{k=0}^{j} B[k][0]+B[j][1]\right).
\]
Esto coincide exactamente con la forma requerida para la siguiente iteración.

Por inducción, la invariante es cierta durante todo el ciclo. Al terminar, con \(i=N\),
se tiene
\[
ans=T(B).
\]
Como \(B\) es óptimo por el teorema de optimización global, el algoritmo retorna el
tiempo mínimo de completación. \(\square\)

\subsection*{Complejidad}

El paso dominante es el ordenamiento del arreglo:
\[
O(N\log N).
\]
El recorrido posterior es lineal:
\[
O(N).
\]
En consecuencia, la complejidad temporal total es
\[
O(N\log N).
\]

Si el ordenamiento se hace in situ, el espacio adicional del recorrido es
\[
O(1),
\]
o bien el correspondiente al algoritmo de ordenamiento usado.

\newpage

\section{Problema 2: Balanced Parentheses}

\subsection*{Resumen del problema}

Se da un arreglo
\[
w[0..n),
\]
tal que para cada \(i\), con \(0\leq i<n\), se cumple que \(w[i]\) es \texttt{(} o
\texttt{)}. El problema tiene dos partes.

En la primera, se debe decidir si la cadena completa está balanceada. En la segunda,
se debe calcular la longitud de una subsecuencia balanceada más larga de la cadena.
En ambos casos se pide una solución en tiempo lineal.

Para la primera parte, basta recorrer la cadena manteniendo el balance entre aperturas
y cierres. Para la segunda parte, la estrategia voraz consiste en recorrer la cadena de
izquierda a derecha y emparejar cada paréntesis de cierre con algún paréntesis de
apertura previamente disponible, siempre que esto sea posible.

\subsection*{Notación auxiliar}

Para cada \(i\), con \(0\leq i\leq n\), defínase:
\[
a(i)=\big|\{j:0\leq j<i \land w[j]=\texttt{(}\}\big|,
\]
\[
c(i)=\big|\{j:0\leq j<i \land w[j]=\texttt{)}\}\big|,
\]
y
\[
b(i)=a(i)-c(i).
\]

Así, \(b(i)\) representa el balance del prefijo \(w[0..i)\).

\subsection*{Parte (a): decidir si una cadena está balanceada}

\paragraph{Caracterización.}
La cadena \(w[0..n)\) está balanceada si y solo si:
\[
b(n)=0
\qquad\text{y}\qquad
\forall i\; (0\leq i\leq n \Rightarrow b(i)\geq 0).
\]

\paragraph{Demostración.}
\textbf{(\(\Rightarrow\))} Se procede por inducción estructural sobre la definición de cadena
balanceada.

Si \(w\) es la cadena vacía, entonces \(b(0)=0\) y la propiedad es inmediata.

Si \(w=(x)\), donde \(x\) es balanceada, entonces todo prefijo interno de \(x\) tiene balance
no negativo; al agregar el par de paréntesis externos, los balances de los prefijos
correspondientes siguen siendo no negativos y el balance total vuelve a \(0\).

Si \(w=xy\), donde \(x\) y \(y\) son balanceadas, entonces todo prefijo de \(x\) tiene balance
no negativo; y todo prefijo que se extiende dentro de \(y\) tiene balance igual a
\[
b(|x|)+b_y(t)=0+b_y(t)\geq 0.
\]
Además,
\[
b(|w|)=b(|x|)+b_y(|y|)=0+0=0.
\]

\textbf{(\(\Leftarrow\))} Se procede por inducción en \(n\).

Si \(n=0\), la cadena vacía es balanceada.

Supóngase \(n>0\). Como \(b(i)\geq 0\) para todo \(i\) y \(b(n)=0\), necesariamente
\[
w[0]=\texttt{(}.
\]
Sea \(r>0\) el menor índice tal que
\[
b(r)=0.
\]
Entonces \(w[r-1]=\texttt{)}\), y puede escribirse
\[
w=(x)y,
\]
donde
\[
x=w[1..r-1)
\qquad\text{y}\qquad
y=w[r..n).
\]
Por minimalidad de \(r\), los prefijos de \(x\) tienen balance no negativo y el balance
total de \(x\) es \(0\). De igual manera, \(y\) también satisface la misma caracterización.
Por hipótesis inductiva, \(x\) y \(y\) son balanceadas. Luego \(w=(x)y\) es balanceada.
\(\square\)

\paragraph{Pseudocódigo.}
\[
\begin{array}{l}
\textsc{IsBalanced}(w[0..n)):\\
\quad bal \gets 0 \\
\\
\quad \textbf{for } i = 0 \textbf{ to } n-1 \textbf{ do}\\
\quad \quad \textbf{if } w[i] = \texttt{(} \textbf{ then}\\
\quad \quad \quad bal \gets bal + 1\\
\quad \quad \textbf{else}\\
\quad \quad \quad bal \gets bal - 1\\
\quad \quad \quad \textbf{if } bal < 0 \textbf{ then return false}\\
\\
\quad \textbf{return } (bal = 0)
\end{array}
\]

\paragraph{Ejemplo.}
Para
\[
w=\texttt{(()())()},
\]
el recorrido de balances es:
\[
\begin{array}{c|cccccccc}
i & 0 & 1 & 2 & 3 & 4 & 5 & 6 & 7\\
\hline
w[i] & \texttt{(} & \texttt{(} & \texttt{)} & \texttt{(} & \texttt{)} & \texttt{)} & \texttt{(} & \texttt{)}\\
b(i+1) & 1 & 2 & 1 & 2 & 1 & 0 & 1 & 0
\end{array}
\]
Nunca aparece un balance negativo y el balance final es \(0\), así que la cadena es
balanceada.

\paragraph{Invariante del ciclo.}
Si el algoritmo no ha retornado \texttt{false} antes de comenzar la iteración \(i\),
entonces:
\[
bal=b(i)
\]
y
\[
\forall j\; (0\leq j\leq i \Rightarrow b(j)\geq 0).
\]

\paragraph{Demostración.}
Para \(i=0\), antes del ciclo,
\[
bal=0=b(0),
\]
y la propiedad sobre los prefijos es inmediata.

Supóngase la invariante cierta al comienzo de la iteración \(i\), con \(0\leq i<n\).

Si \(w[i]=\texttt{(}\), el algoritmo incrementa \(bal\), de modo que
\[
bal=b(i)+1=b(i+1).
\]
Además, como \(b(i)\geq 0\), también
\[
b(i+1)\geq 0.
\]

Si \(w[i]=\texttt{)}\), el algoritmo decrementa \(bal\), y obtiene
\[
bal=b(i)-1=b(i+1).
\]
Si este valor es negativo, retorna \texttt{false}; en efecto, ya existe un prefijo con balance
negativo. Si no es negativo, entonces la propiedad sigue siendo válida para todos los
prefijos hasta \(i+1\).

Por inducción, la invariante es cierta. Si el algoritmo retorna \texttt{true}, entonces
no apareció ningún prefijo con balance negativo y además el balance final es \(0\).
Por la caracterización, la cadena está balanceada. Si retorna \texttt{false}, entonces
apareció un prefijo con balance negativo, y por la misma caracterización la cadena no
puede estar balanceada. \(\square\)

\paragraph{Complejidad.}
El algoritmo recorre la cadena una sola vez y realiza trabajo \(O(1)\) por posición. Su
complejidad temporal es
\[
O(n),
\]
y usa espacio adicional
\[
O(1).
\]

\subsection*{Parte (b): longitud de una subsecuencia balanceada más larga}

\paragraph{Función objetivo.}
Para cada \(i\), con \(0\leq i\leq n\), sea
\[
L(i)=\text{longitud de una subsecuencia balanceada más larga de } w[0..i).
\]

El objetivo del problema es calcular
\[
L(n).
\]

\paragraph{Estrategia voraz.}
Se recorre la cadena de izquierda a derecha. Cada vez que aparece un \texttt{(}, se
considera como una apertura disponible para emparejar en el futuro. Cada vez que
aparece un \texttt{)}, si existe al menos una apertura disponible, se empareja de inmediato
con una de ellas y se incrementa en \(2\) la longitud de la subsecuencia construida.

La apuesta voraz es:

\begin{quote}
Si un paréntesis de cierre puede emparejarse con alguna apertura ya vista, conviene
emparejarlo inmediatamente.
\end{quote}

\paragraph{Pseudocódigo.}
\[
\begin{array}{l}
\textsc{LongestBalancedSubsequence}(w[0..n)):\\
\quad ab \gets 0 \\
\quad ans \gets 0 \\
\\
\quad \textbf{for } i = 0 \textbf{ to } n-1 \textbf{ do}\\
\quad \quad \textbf{if } w[i] = \texttt{(} \textbf{ then}\\
\quad \quad \quad ab \gets ab + 1\\
\quad \quad \textbf{else if } ab > 0 \textbf{ then}\\
\quad \quad \quad ab \gets ab - 1\\
\quad \quad \quad ans \gets ans + 2\\
\\
\quad \textbf{return } ans
\end{array}
\]

\paragraph{Ejemplo.}
Considérese
\[
w=\texttt{(()))(())}.
\]
El algoritmo produce el siguiente recorrido:

\[
\begin{array}{c|ccccccccc}
i & 0 & 1 & 2 & 3 & 4 & 5 & 6 & 7 & 8\\
\hline
w[i] & \texttt{(} & \texttt{(} & \texttt{)} & \texttt{)} & \texttt{)} & \texttt{(} & \texttt{(} & \texttt{)} & \texttt{)}\\
ab & 1 & 2 & 1 & 0 & 0 & 1 & 2 & 1 & 0\\
ans & 0 & 0 & 2 & 4 & 4 & 4 & 4 & 6 & 8
\end{array}
\]

Por tanto, la longitud de una subsecuencia balanceada más larga es
\[
8.
\]

\paragraph{Teorema de optimización local.}
Sea \(i\), con \(0\leq i<n\).

\begin{enumerate}[label=\alph*)]
    \item Si \(w[i]=\texttt{(}\), entonces
    \[
    L(i+1)=L(i).
    \]

    \item Si \(w[i]=\texttt{)}\) y
    \[
    a(i)=\frac{L(i)}{2},
    \]
    entonces
    \[
    L(i+1)=L(i).
    \]

    \item Si \(w[i]=\texttt{)}\) y
    \[
    a(i)>\frac{L(i)}{2},
    \]
    entonces
    \[
    L(i+1)=L(i)+2.
    \]
\end{enumerate}

\paragraph{Demostración.}
\textbf{a)} Si el símbolo nuevo es \texttt{(}, ninguna subsecuencia balanceada de
\(w[0..i+1)\) puede aumentar su longitud usando únicamente ese símbolo al final, pues
un paréntesis de apertura sin su cierre correspondiente no aporta una subsecuencia
balanceada nueva. Luego
\[
L(i+1)=L(i).
\]

\textbf{b)} Si el símbolo nuevo es \texttt{)} y además
\[
a(i)=\frac{L(i)}{2},
\]
entonces cualquier subsecuencia balanceada de \(w[0..i+1)\) usa a lo sumo \(a(i)\)
aperturas provenientes del prefijo \(w[0..i)\). Por tanto, su longitud es a lo sumo
\[
2a(i)=L(i).
\]
Como ya existe una subsecuencia balanceada de longitud \(L(i)\) en el prefijo anterior,
se concluye que
\[
L(i+1)=L(i).
\]

\textbf{c)} Si el símbolo nuevo es \texttt{)} y además
\[
a(i)>\frac{L(i)}{2},
\]
entonces hay más aperturas en el prefijo \(w[0..i)\) que las que usa una subsecuencia
balanceada óptima de longitud \(L(i)\). En particular, existe al menos una apertura
disponible que puede emparejarse con el símbolo nuevo, así que se puede construir una
subsecuencia balanceada de longitud \(L(i)+2\). Por tanto,
\[
L(i+1)\geq L(i)+2.
\]

Por otra parte, cualquier subsecuencia balanceada de \(w[0..i+1)\) o bien no usa el
símbolo \(w[i]\), en cuyo caso su longitud es a lo sumo \(L(i)\), o bien sí lo usa. En
este último caso, al eliminar de ella el par formado por \(w[i]\) y su apertura
correspondiente, se obtiene una subsecuencia balanceada de \(w[0..i)\), cuya longitud
es a lo sumo \(L(i)\). Por tanto, la longitud original era a lo sumo \(L(i)+2\).

De las dos desigualdades se concluye:
\[
L(i+1)=L(i)+2.
\]
\(\square\)

\paragraph{Invariante del ciclo.}
Para cada \(i\), con \(0\leq i\leq n\), después de procesar el prefijo \(w[0..i)\), se cumple:
\[
ans=L(i)
\]
y
\[
ab=a(i)-\frac{L(i)}{2}.
\]

\paragraph{Demostración.}
Se procede por inducción sobre \(i\).

\textbf{Caso base:} \(i=0\). Antes del ciclo,
\[
ans=0
\qquad\text{y}\qquad
ab=0.
\]
Además,
\[
L(0)=0
\qquad\text{y}\qquad
a(0)=0.
\]
Luego
\[
ans=L(0)
\qquad\text{y}\qquad
ab=a(0)-\frac{L(0)}{2}.
\]

\textbf{Paso inductivo:} supóngase la propiedad cierta para algún \(i\), con \(0\leq i<n\).

Si \(w[i]=\texttt{(}\), el algoritmo incrementa \(ab\) y deja \(ans\) intacto. Por el
teorema de optimización local,
\[
L(i+1)=L(i).
\]
Además,
\[
a(i+1)=a(i)+1.
\]
Luego,
\[
ans=L(i)=L(i+1)
\]
y
\[
ab+1
=
a(i)-\frac{L(i)}{2}+1
=
a(i+1)-\frac{L(i+1)}{2}.
\]

Si \(w[i]=\texttt{)}\) y \(ab=0\), entonces por hipótesis inductiva
\[
a(i)=\frac{L(i)}{2}.
\]
Por el teorema de optimización local,
\[
L(i+1)=L(i).
\]
El algoritmo no modifica ni \(ab\) ni \(ans\), así que
\[
ans=L(i+1)
\]
y, como \(a(i+1)=a(i)\),
\[
ab=0=a(i+1)-\frac{L(i+1)}{2}.
\]

Si \(w[i]=\texttt{)}\) y \(ab>0\), entonces por hipótesis inductiva
\[
a(i)-\frac{L(i)}{2}>0.
\]
Por el teorema de optimización local,
\[
L(i+1)=L(i)+2.
\]
El algoritmo hace
\[
ab \gets ab-1
\qquad\text{y}\qquad
ans \gets ans+2.
\]
Por tanto,
\[
ans+2=L(i)+2=L(i+1),
\]
y como \(a(i+1)=a(i)\),
\[
ab-1
=
a(i)-\frac{L(i)}{2}-1
=
a(i+1)-\frac{L(i+1)}{2}.
\]

En todos los casos la invariante se preserva. Por inducción, es cierta para todo \(i\).
Al terminar el ciclo, con \(i=n\), se tiene
\[
ans=L(n),
\]
que es exactamente la longitud de una subsecuencia balanceada más larga. \(\square\)

\paragraph{Teorema de optimización global.}
La invocación de
\[
\textsc{LongestBalancedSubsequence}(w)
\]
retorna la longitud de una subsecuencia balanceada más larga de \(w[0..n)\).

\paragraph{Demostración.}
Es una consecuencia inmediata de la invariante del ciclo para \(i=n\). \(\square\)

\paragraph{Complejidad.}
El algoritmo recorre la cadena una sola vez, realizando trabajo constante en cada
posición. Por tanto, su complejidad temporal es
\[
O(n),
\]
y usa espacio adicional
\[
O(1).
\]

\section*{Conclusión}

En el problema de El Goog, la estructura del sistema permite un diseño voraz clásico:
escoger primero los trabajos con mayor tiempo de finalización en PC. La correctitud
se demuestra mediante un teorema de optimización local por intercambio y un teorema
de optimización global que fija el prefijo óptimo de manera inductiva.

En Balanced Parentheses, la primera parte se resuelve con una verificación lineal del
balance de prefijos, mientras que la segunda admite una estrategia voraz lineal que
empareja cada cierre con una apertura disponible tan pronto como ello es posible.
En ambos casos, las invariantes explícitas justifican formalmente la correctitud del
recorrido iterativo.

\end{document}